% archive.tex — InnovationZero archive: every complete text the paper compares, byte for byte.
% The entries are generated into the build directory (archive-index.tex) by scripts/check.sh
% --prepare, from evidence/artifact-manifest.json and evidence/pilot-2026-10-22/programs.tsv, after
% every byte count and SHA-256 has been verified. Nothing here is typed by hand.
\documentclass[10pt]{article}
\usepackage{iz}
\externaldocument[M-]{InnovationZero}
\renewcommand\thesection{\arabic{section}}
\renewcommand\IZShortTitle{InnovationZero \textperiodcentered{} Archive}
\renewcommand\IZTitle{Archive}
\renewcommand\IZSubtitle{The General Regressor Challenge, the four frontier entrants and the ten
  engine programs compared in \emph{InnovationZero}, printed in full}
\renewcommand\IZAuthors{Jepson Taylor\quad Jordan Morrow\quad Alton Alexander}
\renewcommand\IZDate{22 October 2026 \textperiodcentered{} repository \code{neverhuman/JopeDime}}
\renewcommand\IZStatus{SUPPLEMENTARY MATERIAL \textperiodcentered{} ARCHIVE}
\hypersetup{pdftitle={InnovationZero: Archive}, pdfauthor={Jepson Taylor, Jordan Morrow, Alton Alexander}}

\DefineVerbatimEnvironment{ArchiveText}{Verbatim}{}
\newcommand\archivepart[2]{\clearpage
  \phantomsection\addcontentsline{toc}{part}{#1}%
  {\noindent\sffamily\sbseries\Large\color{izaccent}#1\par}\medskip
  \noindent #2\par}
% #1 kind (text|rust), #2 title, #3 id, #4 path, #5 bytes, #6 lines, #7 SHA-256, #8 snapshot
\newcommand\archiveentry[8]{\clearpage
  \section{#2}\label{arc:#3}
  \begin{center}\tablesf
  \begin{tabularx}{\linewidth}{@{}lL@{}}
  \toprule
  File & \nolinkurl{docs/paper/innovationzero/#4} \\
  Size & #5 bytes, #6 lines \\
  SHA-256 & \nolinkurl{#7} \\
  Snapshot & #8 \\
  \bottomrule
  \end{tabularx}
  \end{center}
  \ifx\relax#1\relax\else\csname archive@#1\endcsname{#4}\fi}
\expandafter\def\csname archive@rust\endcsname#1{\lstinputlisting[language=Rust]{#1}}
\expandafter\def\csname archive@text\endcsname#1{\VerbatimInput[fontsize=\fontsize{7.2pt}{8.6pt},
  breaklines=true, breakanywhere=true, breaksymbolleft={\textcolor{izrule}{$\hookrightarrow$}},
  breaksymbolright={}, numbers=left, numbersep=7pt,
  formatcom=\color{izink}\ttfamily, tabsize=4, xleftmargin=16pt]{#1}}

\begin{document}
\izmaketitle
\noindent This volume prints, in full and without edits, every text whose content the paper analyses:
the challenge as its entrants received it and as its organizer wrote it, the four frozen frontier
entrants, and the ten programs the engine wrote for the pilot datasets. Each entry is typeset
directly from the archived file named in its header, after the build has checked that file's byte
count and SHA-256 against the manifest; snapshot commits identify where in the repository's history
each archived copy was taken. Line numbers are the file's own; a hooked arrow marks a line wrapped
for the page. Wrapping and page breaks are presentation only: the files and their hashes are the
authority. Printing a source establishes what was compiled and scored; it does not establish how
fairly a session was run or how well a program generalizes, which the paper discusses with the
receipts.

{\small\tableofcontents}
\archivepart{The challenge and the four frontier entrants}{The participant packet and master file of
the General Regressor Challenge (Section~\ref{M-sec:challenge} of the paper), followed by the four
frontier-written regressors in order of their Elo on the 64-family panel.}
\input{archive-index}
\end{document}
